\section{Methods}
\label{sec:methods}

The analysis assumes a linear measurement model in which an observed response is
a weighted combination of two inputs plus an additive disturbance. The weights
and the noise scale are fixed in advance and collected in \cref{tab:params} so
that the results section can refer to concrete values rather than abstract
symbols. Holding these parameters constant lets the closed-form result in
\cref{sec:results} be stated without further qualification.

\begin{table}[htbp]
\centering
\caption{Fixed parameters of the measurement model.}
\label{tab:params}
\begin{tabular}{lcl}
\toprule
Symbol & Value & Interpretation \\
\midrule
$\beta_1$    & $0.60$ & weight on the first input \\
$\beta_2$    & $0.25$ & weight on the second input \\
$\sigma$     & $0.10$ & standard deviation of the disturbance \\
$n$          & $500$  & number of observations \\
\bottomrule
\end{tabular}
\end{table}

Estimation proceeds by ordinary least squares applied to the two inputs jointly.
Because the design matrix is fixed and the disturbance is homoskedastic with the
scale given in \cref{tab:params}, the estimator is unbiased and its sampling
variance has a particularly simple form, which the next section records
explicitly.
